\chapter{Transferencia de calor y refrigeración}\label{ch:17}
\objetivos{Cerrar un balance térmico de pared y refrigerante; distinguir temperatura de gas, recuperación y material; interpretar resistencias, transitorios y recuperación regenerativa sin confundir energía conservada con reversibilidad.}
\section{La carga térmica y sus temperaturas de referencia}
Un flujo caliente puede transferir energía por convección y radiación. La temperatura que impulsa la convección compresible junto a una pared no es necesariamente la temperatura estática de la corriente exterior. Se utiliza una temperatura de recuperación o pared adiabática $T_{aw}$, definida por el estado que alcanzaría una pared sin intercambio neto bajo esas condiciones.

Una representación de gas caloríficamente perfecto es
\[
 T_{aw}=T+ r_h\frac{u^2}{2c_p},
\]
donde $r_h$ es un factor de recuperación ligado al transporte en la capa límite. No se identifica automáticamente con uno ni con un rendimiento propulsivo. El flujo convectivo se modela como $q''_g=h_g(T_{aw}-T_{wg})$. El coeficiente $h_g$ depende del campo y propiedades; no es una constante universal de la cámara.\cite{MIT07,MIT09}
\section{Pared plana y resistencias térmicas}
En una pared plana de espesor $\delta$, conductividad $k$ constante, régimen estacionario y conducción unidimensional sin generación,
\[
 q''=k\frac{T_{wg}-T_{wl}}{\delta}.
\]
Si el lado refrigerante tiene $q''=h_l(T_{wl}-T_l)$, sumar caídas de temperatura da
\begin{equation}\label{eq:thermalresistance}
 q''=\frac{T_{aw}-T_l}{1/h_g+\delta/k+1/h_l}.
\end{equation}
Las resistencias por unidad de área tienen unidades $\mathrm{m^2K/W}$. Las temperaturas de pared se recuperan restando cada salto, no asignando al metal la temperatura del gas.

La igualdad de los tres flujos exige estacionariedad y ausencia de otras ramas de transferencia. En una pared transitoria parte de la energía se almacena; con radiación aparece otra contribución. Para geometría cilíndrica o áreas distintas, las resistencias deben referirse a áreas coherentes y la conducción contiene logaritmos radiales. Usar directamente la fórmula plana requiere justificar espesor pequeño respecto al radio.
\section{Capacidad térmica del refrigerante}
Un circuito que absorbe potencia $\dot Q_c$ satisface
\[
 \dot Q_c=\dot m_c(h_{out}-h_{in})
\]
si se desprecia trabajo, cambio cinético y altura. Para $c_{p,c}$ aproximadamente constante y una sola fase,
\[
 \Delta T_c=\frac{\dot Q_c}{\dot m_cc_{p,c}}.
\]
Esta ecuación determina un aumento medio de entalpía; no garantiza que la pared local permanezca admisible. Una potencia total moderada puede concentrarse en una región pequeña con flujo superficial elevado.

Si hay cambio de fase, se debe usar entalpía y fracción de vapor en lugar de extender $c_p\Delta T$ a través de la transición. Las propiedades y la caída de presión pueden cambiar intensamente. El estado del refrigerante se comprueba localmente antes de asumir que conserva densidad y capacidad calorífica constantes.
\section{Refrigeración regenerativa y segunda ley}
En un sistema regenerativo, el propulsante absorbe calor de pared y después entra en la cámara. Para una frontera que encierra gas, pared y circuito hasta esa entrada, gran parte del calor intercambiado es interno. La energía recuperada puede volver al proceso propulsivo. Eso no vuelve reversible la transferencia: mover calor desde gas muy caliente a fluido más frío genera entropía.

Para un intercambio diferencial entre temperaturas $T_h>T_c$,
\[
 dS_{gen}=\delta Q\left(\frac1{T_c}-\frac1{T_h}\right)>0.
\]
Además existen pérdidas de presión y potencia de bombeo. Decir que se recupera energía no implica que se recupere íntegramente su capacidad de producir trabajo. El análisis de ciclo debe incluir ambas cosas.\cite{MIT09}
\section{Película, ablación, acumulación y radiación}
Una película modifica la composición y temperatura de la región próxima a pared. Puede disminuir carga térmica, pero también cambiar mezcla y prestación. El refrigerante consumido forma parte del balance de masa total; excluirlo de $I_{sp}$ genera una comparación inconsistente.

Una protección ablativa utiliza transformaciones y pérdida de material. Un balance superficial idealizado puede escribirse
\[
 q''_{in}=q''_{cond}+\dot m''_a h_{ef}+q''_{rad},
\]
con $h_{ef}$ definido para agrupar procesos del modelo. No debe interpretarse como una constante universal sin datos de material y condiciones. La frontera se mueve y puede modificar geometría.

Una estructura absorbente almacena energía sensible: $mc\,dT/dt=\dot Q_{in}-\dot Q_{out}$ en un modelo concentrado. La radiación neta hacia un entorno representativo puede aproximarse por $\epsilon\sigma(T_w^4-T_{ent}^4)$ bajo hipótesis de superficies y factores de vista adecuados. Estos mecanismos pueden coexistir; se suman mediante un balance, no escogiendo arbitrariamente la temperatura más favorable.\cite{MIT10}
\section{Escalas y acoplamiento estructural}
El número de Biot $Bi=h\delta/k$ compara resistencia interna de conducción con resistencia superficial. Si es pequeño y la geometría permite una longitud característica adecuada, un modelo térmico concentrado puede ser razonable. El número de Fourier $Fo=\alpha t/\delta^2$ compara duración con difusión interna. Son indicadores del modelo, no valores admisibles de material.

Un gradiente térmico produce dilataciones incompatibles si la estructura está restringida. Una escala de tensión térmica elástica es $E\alpha_T\Delta T$, modificada por restricción y coeficiente de Poisson. Combinarla con presión requiere geometría, compatibilidad y una ley constitutiva. La temperatura también puede cambiar resistencia, fluencia y fatiga. Un balance de calor correcto no basta para demostrar integridad estructural.\cite{MIT06}
\figura{c17}{Red térmica ideal. Las temperaturas de pared se obtienen de las caídas parciales; no se igualan a la temperatura del gas.}
\section{Caso desarrollado}
\begin{studiumworked}{M17. Pared y refrigerante: cierre térmico · INTERMEDIATE}
\textbf{Enunciado original.} Una pared plana académica tiene espesor 1.50 mm y $k=18$ W/(m K). El gas tiene $T_{aw}=1134$ K y $h_g=1970$ W/(m$^2$ K); el refrigerante, 290 K y $h_l=5000$ W/(m$^2$ K). Para área 0.012 m$^2$, caudal 0.15 kg/s y $c_p=2200$ J/(kg K), obtenga flujo, temperaturas de caras y aumento térmico del refrigerante.

\textbf{Datos utilizados.} Convección gas: $1970\,\mathrm{W\,m^{-2}\,K^{-1}}$; Convección líquido: $5000\,\mathrm{W\,m^{-2}\,K^{-1}}$; Espesor: $0.0015\,\mathrm{m}$; Conductividad: $18\,\mathrm{W\,m^{-1}\,K^{-1}}$; Recuperación: $1134\,\mathrm{K}$; Refrigerante: $290\,\mathrm{K}$.

\textbf{Incógnitas.} q, temperaturas de pared y potencia.

\textbf{Hipótesis.} Estacionario unidimensional, propiedades constantes. Sin radiación ni almacenamiento; temperatura de refrigerante representativa.

\textbf{Modelo físico.} Resistencias térmicas y balance de entalpía.
\[q^{\prime\prime}=(T_{aw}-T_l)/(1/h_g+\delta/k+1/h_l)\]
\[T_{wg}=T_{aw}-q^{\prime\prime}/h_g,\quad T_{wl}=T_l+q^{\prime\prime}/h_l\]
\[\dot Q=q^{\prime\prime}A,\quad \Delta T_l=\dot Q/(\dot m_lc_{p,l})\]
\textbf{Despejes y sustituciones.}
\begin{enumerate}
\item Las resistencias suman $7.909475\times10^{-4}\,\mathrm{m^2K/W}$. Dividir el salto $844\,\mathrm{K}$ da $1.067075\times10^{6}\,\mathrm{W/m^2}$.
\item La caída del lado gas es $541.6622\,\mathrm{K}$, la de pared $88.92288\,\mathrm{K}$ y la del refrigerante $213.4149\,\mathrm{K}$. Así las caras quedan a $592.3378\,\mathrm{K}$ y $503.4149\,\mathrm{K}$.
\item Multiplicar por área da $12804.89\,\mathrm{W}$. El balance del refrigerante produce $\Delta T=12804.89/(0.15\cdot2200)=38.80271$ K. Esa subida pequeña se compara con el salto gas-líquido para valorar la temperatura representativa inicial.
\end{enumerate}
\textbf{Resultados.} $q^{\prime\prime}=1.067075\times10^{6}$ W/m$^2$; $T_{wg}=592.3378$ K; $T_{wl}=503.4149$ K; $\Delta T_l=38.80271$ K.

\textbf{Interpretación.} Las temperaturas metálicas se calculan entre las del gas y líquido. La potencia total y la temperatura local de pared son restricciones distintas.

\textbf{Comprobación.} $h_g(T_{aw}-T_{wg})=k(T_{wg}-T_{wl})/\delta=h_l(T_{wl}-T_l)=1.067075\times10^{6}$ W/m$^2$. El calor absorbido coincide con $\dot m c_p\Delta T$.

\textbf{Límites.} No se comprueban ebullición, resistencia del material ni geometría de un motor; si la subida del líquido fuese grande, habría que resolver la variación axial.
\end{studiumworked}

\section{Ejercicios propuestos}
\begin{enumerate}
\item Deduzca \eqref{eq:thermalresistance} y recupere las temperaturas de ambas caras de pared.
\item Para una potencia absorbida de 24 kW, caudal académico $0.80\,\mathrm{kg/s}$ y $c_p=2500\,\mathrm{J/(kg\,K)}$, calcule el aumento de temperatura y cierre energía.
\item Explique por qué recuperar el calor en el propulsante no elimina generación de entropía.
\item Distinga qué balance cambia si el refrigerante hierve, si la pared acumula energía o si una superficie se abla.
\end{enumerate}
\resumen{La temperatura del material se obtiene de un problema de transferencia, no de igualarla a la del gas. La capacidad global del refrigerante y el flujo local son restricciones distintas. La regeneración recupera energía dentro del sistema y aun así genera entropía y pérdidas de presión.}


\studiumchaptercheck
